\lecture{14}{20. Oktober 2025}{Torsional Deformation}

\begin{problemwithimage}{f10_1}{10-1}
  The tube is subjected to a torque of \qty{750}{\N.\m}.
  Determine the amount of this torque that is resisted by the gray shaded section.
  Solve the problem two ways: (a) by using the torsion formula, (b) by finding the resultant of the shear stress distribution.
\end{problemwithimage}

\begin{derivation}
  \begin{enumerate}[label=(\alph*)]
    \item The torsion formula is
      \begin{equation*}
        \tau(r) = \frac{T r}{J}
      .\end{equation*}
      With $\tau$ being the shear stress at distance $r$ from the center, $T$ the internal torque, $J$ the polar moment of inertia.

      The polar moment of inertia of the rod is
      \begin{equation*}
        J = \frac{\pi}{2} \left( r_o^{4} - r_i^{4} \right)
      .\end{equation*}
      
      The share of the torque carried by the shaded part of the rod will be proportional to its share of the entire polar moment of inertia:
      \begin{equation*}
        \frac{T_{\mathrm{gray}}}{J_{\mathrm{gray}}} = \frac{T_{\mathrm{tot}}}{J_{\mathrm{tot}}} \implies T_{\mathrm{gray}} = T_{\mathrm{tot}} \frac{J_{\mathrm{gray}}}{J_{\mathrm{tot}}}
      .\end{equation*}
      So,
      \begin{equation*}
      T_{\mathrm{gray}} = \qty{750}{\N.\m} \cdot \frac{\frac{\pi}{2} \cdot \left( \left( \qty{100}{\mm}  \right)^4 - \left( \qty{75}{\mm} \right)^4 \right)}{\frac{\pi}{2} \cdot \left( \qty{100}{\mm}  \right)^{4} - \left( \qty{25}{\mm}  \right)^{4}} = \qty{513,9}{\N.\m}
    .\end{equation*}
  \item The torque contributed by an area element is
    \begin{equation*}
      \mathrm{d}T = \tau r \, \mathrm{d}A
    \end{equation*}
    so with
    \begin{equation*}
      \tau(r) = \frac{T r}{J_{\mathrm{tot}}}
    \end{equation*}
    we get
    \begin{align*}
      \mathrm{d}T &= \frac{T r^2}{J_{\mathrm{tot}}} \, \mathrm{d}A \\
      T_{\mathrm{gray}} &= \frac{T}{J_{\mathrm{tot}}} \int_A r^2 \, \mathrm{d}A \\
      T_{\mathrm{gray}} &= \frac{T}{J_{\mathrm{tot}}} J_{\mathrm{gray}} \\
      &= \qty{513,9}{\N.\m} 
    .\end{align*}
  \end{enumerate}
\end{derivation}

\finalresult{T_{\mathrm{gray}} = \qty{513,9}{\N.\m}}

\begin{problemwithimage}{f10_26}{10-26}
  The \qty{160}{\mm} diameter solid shaft is supported by bearings at $A$ and $E$.
  Determine the maximum shear stress developed in each segment of the shaft.
\end{problemwithimage}

\begin{derivation}
  The internal torque at a section is the sum of the torques applied to the left of the section.
  So
  \begin{itemize}
    \item $AB$: \qty{0}{\kN.\m}
    \item $BC$: \qty{-40}{\kN.\m} 
    \item $CD$: \qty{20}{\kN.\m} 
    \item $DE$: \qty{50}{\kN.\m}
    \item $EF$: \qty{50}{\kN.\m}
    \item $F+$: \qty{0}{\kN.\m} 
  \end{itemize}
  From the torsion formula we get
  \begin{equation*}
    \tau_{\mathrm{max}} = \frac{Tc}{J} = \frac{16}{\pi d^3} T = \frac{16}{\pi \left( \qty{160}{\mm}  \right)^3} T = \qty{1,24e-6}{\per\mm^3} \cdot T 
  .\end{equation*}
  So for each segment:
\end{derivation}
\finalresult{
\begin{aligned}
    &AB: & \tau_{\mathrm{max}} &= \qty{1.24e-6}{\mm^3} \cdot \qty{0}{\kN.\m} & &= \qty{0}{\MPa}  \\
    &BC: & \tau_{\mathrm{max}} &= \qty{1.24e-6}{\mm^3} \cdot \qty{-40}{\kN.\m} & &= \qty{-49,6}{\MPa}  \\
    &CD: & \tau_{\mathrm{max}} &= \qty{1.24e-6}{\mm^3} \cdot \qty{20}{\kN.\m} & &= \qty{24,8}{\MPa}  \\
    &DE: & \tau_{\mathrm{max}} &= \qty{1.24e-6}{\mm^3} \cdot \qty{50}{\kN.\m} & &= \qty{62,0}{\MPa}  \\
    &EF: & \tau_{\mathrm{max}} &= \qty{1.24e-6}{\mm^3} \cdot \qty{50}{\kN.\m} & &= \qty{62,0}{\MPa}
.\end{aligned}
}
